% Coverage: physical pp. 105--114 (printed pp. 82--91), from the Section 17.2 heading through the paragraph ending immediately before Section 17.3; Eqs. (17.2.1)--(17.2.28), all unnumbered displays, five footnotes, and the centered three-asterisk divider included.
% Figures: none. Tables: none. Uncertainties: none after full-resolution rendered-page review.
\subsection{Renormalization: Direct Analysis}

The simplest non-Abelian gauge theories are renormalizable in the `Dyson' sense that the operators in the Lagrangian density all have dimensionality (in powers of mass) four or less. As we saw in Chapter 12, this guarantees that the infinities in the quantum effective action only appear in terms that could be cancelled by the counterterms in interactions of dimensionality four or less. But there is more to renormalizability than this. The Lagrangian density is constrained by gauge invariance and other symmetries. For a theory to be renormalizable, it is necessary that the infinities in the quantum effective action satisfy the same constraints, up to possible renormalizations of the fields.

The effective action $\Gamma[\chi,K]$ is a complicated functional of both $\chi$ and $K$, about which the symmetry condition \eqref{eq:17.1.9} says complicated things, but fortunately matters are much simpler for the infinite terms in $\Gamma$. We will write the action $S[\chi,K]=I[\chi]+\int d^4x\,\Delta^nK_n$ as the sum of a term $S_R[\chi,K]$ in which masses and coupling constants are set equal to their renormalized values, plus a correction $S_\infty[\chi,K]$, which contains the counterterms that we intend to cancel the infinities from loop graphs. Both $S_R$ and $S_\infty$ must be taken to have the symmetries of the original action $S[\chi,K]$, so the question is whether the infinite parts of the higher-order contributions to $\Gamma$ share the same symmetries, so that they can be cancelled by the counterterms in $S_\infty$.

We may expand $\Gamma$ in a series of terms $\Gamma_N$ arising from diagrams with just $N$ loops, plus contributions from graphs with $N-M$ loops (where $1\leq M\leq N$) involving various counterterms in $S_\infty[\chi,K]$ that will be used to cancel infinities in graphs with a total of $M$ loops:
\begin{equation}
\Gamma[\chi,K]=\sum_{N=0}^{\infty}\Gamma_N[\chi,K].
\tag{17.2.1}\label{eq:17.2.1}
\end{equation}
The symmetry condition \eqref{eq:17.1.10} then reads,\footnote{Such order-by-order relations may be derived formally by repeating the reasoning of Section 16.1. When the action $S_R$ is replaced with $g^{-1}S_R$, the contribution of a connected $L$-loop graph with $I$ internal lines and $V$ vertices is multiplied by a factor $g^{I-V}=g^{L-1}$. If the counterterms in $S_\infty$ associated with $N$-loop diagrams are also provided with factors $g^N$, then $\Gamma_L$ is the value for $g=1$ of the term in $\Gamma$ of order $g^{L-1}$. Eq.~\eqref{eq:17.2.2} then follows by requiring Eq.~\eqref{eq:17.1.10} to hold in each order in $g$. In cgs units the action has the same dimensions as $\hbar$, and so appears in the path integral multiplied with a factor $1/\hbar$, so we can also use $\hbar$ as a loop-counting parameter in place of $g$.} for each $N$,
\begin{equation}
\sum_{N'=0}^{N}\bigl(\Gamma_{N'},\Gamma_{N-N'}\bigr)=0.
\tag{17.2.2}\label{eq:17.2.2}
\end{equation}
In the sum \eqref{eq:17.2.1} the leading term is just $\Gamma_0[\chi,K]=S_R[\chi,K]$, which of course is finite. Suppose that, for all $M<N-1$, all infinities arising from $M$-loop graphs have been cancelled by counterterms in $S_\infty$. Then infinities can appear in Eq.~\eqref{eq:17.2.2} only in the $N'=0$ and $N'=N$ terms, which are equal, and the infinite part of this condition tells us that the infinite part $\Gamma_{N,\infty}$ of $\Gamma_N$ is subject to the condition that
\begin{equation}
\bigl(S_R,\Gamma_{N,\infty}\bigr)=0.
\tag{17.2.3}\label{eq:17.2.3}
\end{equation}
This is a symmetry principle generated by $S_R$, just like that described by Eqs.~\eqref{eq:15.9.16} and \eqref{eq:15.9.17}. Note in particular that the transformation $X\mapsto(S_R,X)$ acts on the external fields $K_n$ as well as the fields $\chi^n$.

Up to this point, we have used none of the special properties of a renormalizable Yang--Mills theory. Now note that, according to the general power-counting rules of renormalization theory, with all infinities cancelled in subgraphs of $\Gamma_N$, the infinite part $\Gamma_{N,\infty}[\chi,K]$ of $\Gamma_N[\chi,K]$ can only be a sum of products of fields (including $K$) and their derivatives of dimensionality (in powers of mass) four or less. Finally, the arguments of Section 16.4 show that $\Gamma[\chi,K]$ and hence also $\Gamma_{N,\infty}[\chi,K]$ are invariant under all of the \emph{linearly realized} symmetry transformations under which the action is invariant. (As described below, these are: Lorentz transformations, \emph{global} gauge transformations, antighost translations, and the ghost phase transformations associated with ghost number conservation. Of course, the auxiliary fields $K_n$ must be assigned suitable transformation properties under these symmetry transformations.) These conditions together with Eq.~\eqref{eq:17.2.3} will suffice to tell us all we need to know about the structure of $\Gamma_{N,\infty}[\chi,K]$.

To implement these conditions, we need to know the dimensionalities of the external fields $K_n$. If a field $\chi^n$ has dimensionality $d_n$ (that is, $d_n$ powers of mass), then $\Delta^n$ correspondingly has dimensionality $d_n+1$ (as can be seen by inspection of the BRST transformation rules \eqref{eq:15.7.7}--\eqref{eq:15.7.11}), so in order for $\int d^4x\,K_n\Delta^n$ to be dimensionless, $K_n$ must have dimensionality $3-d_n$. The fields $A^{\alpha\mu}$, $\omega^\alpha$, and $\omega^{\alpha *}$ all have dimensionalities $d_n=+1$, so the corresponding $K_n$ all have dimensionalities $+2$. (We do not introduce any external field corresponding to $h^\alpha$, because this field is BRST-invariant.) Any spin $1/2$ matter fields $\psi_\ell$ have dimensionalities $3/2$, and the corresponding $K_n$ thus also have dimensionalities $3/2$. Thus a dimension four quantity like $\Gamma_{N,\infty}[\chi,K]$ is at most quadratic in any of the $K_n$. Furthermore terms that are of second order in the $K_n$ cannot involve any other fields, except that a term of second order in the $K_n$ for spin $1/2$ matter fields may involve at most one additional field of dimensionality unity.

We can now use ghost number conservation to show that in fact $\Gamma_{N,\infty}[\chi,K]$ does not contain any terms of second order in the $K_n$. For this purpose, we also need the ghost quantum numbers of the $K_n$. If $\chi^n$ has ghost quantum number $\gamma_n$, then $\Delta^n$ has ghost quantum number $\gamma_n+1$, so $K_n$ must be assigned ghost quantum number $-\gamma_n-1$. The ghost quantum numbers of the fields $A^{\alpha\mu}$, $\psi^\ell$, $\omega^\alpha$, and $\omega^{\alpha *}$ are respectively $0$, $0$, $+1$, and $-1$, so the corresponding external fields $K_n$ have ghost quantum numbers $-1$, $-1$, $-2$, and $0$, respectively. This rules out any terms in $\Gamma_{N,\infty}[\chi,K]$ that are of second order in the $K_n$, with the one possible exception of a term of second order in the external fields $K_\alpha^*$ associated with $\omega^{\alpha *}$ (and involving no other fields). However, these last terms are also forbidden, for a different reason. The BRST transformation of $\omega^{\alpha *}$ is linear in the fields, with
\begin{equation}
\Delta^{\alpha *}=-h^\alpha,
\tag{17.2.4}\label{eq:17.2.4}
\end{equation}
so here
\[
\frac{\delta_L\Gamma_{N,\infty}[\chi,K]}{\delta K_\alpha^*}
=\langle\Delta^{\alpha *}\rangle_{J_\chi,K}=-h^\alpha
\]
is independent of $K_\alpha^*$. It follows that $\Gamma_{N,\infty}[\chi,K]$ is linear in $K^{\alpha,*}$, and depends on $K^{\alpha,*}$ only through a term $-\int d^4x\,K_\alpha^*h^\alpha$. (Both $K_\alpha^*$ and $h^\alpha$ are bosonic, so their order is immaterial.) In particular, for $N>0$, $\Gamma_{N,\infty}[\chi,K]$ is independent of $K^{\alpha*}$.

We have seen that $\Gamma_{N,\infty}[\chi,K]$ is at most linear in all of the $K_n$. We will write it as
\begin{equation}
\Gamma_{N,\infty}[\chi,K]=\Gamma_{N,\infty}[\chi,0]
+\int d^4x\,\mathcal D_N^n[\chi;x]K_n(x).
\tag{17.2.5}\label{eq:17.2.5}
\end{equation}
We also recall that $S_R$ has the $K$ dependence:
\[
S_R[\chi,K]=S_R[\chi]+\int d^4x\,\Delta^n[\chi;x]K_n(x).
\]
The terms in \eqref{eq:17.2.2} of zeroth and first order in $K$ therefore give\footnote{The second terms in Eqs.~\eqref{eq:17.2.6} and \eqref{eq:17.2.7} have been put into the form shown here by recalling that, because $\chi^n$ and $K_n$ have opposite statistics, for any bosonic functionals $A$ and $B$, \[ \frac{\delta_RA}{\delta\chi^n}\frac{\delta_LB}{\delta K_n} =-\frac{\delta_LA}{\delta\chi^n}\frac{\delta_RB}{\delta K_n} =-\frac{\delta_RB}{\delta K_n}\frac{\delta_LA}{\delta\chi^n}. \]}
\begin{equation}
\int d^4x\left[
\Delta^n[\chi;x]\frac{\delta_L\Gamma_{N,\infty}[\chi,0]}{\delta\chi^n(x)}
+\mathcal D_N^n[\chi;x]\frac{\delta_LS_R[\chi]}{\delta\chi^n(x)}
\right]=0
\tag{17.2.6}\label{eq:17.2.6}
\end{equation}
and
\begin{equation}
\int d^4x\left[
\Delta^n[\chi;x]\frac{\delta_L\mathcal D_N^m[\chi;y]}{\delta\chi^n(x)}
+\mathcal D_N^n[\chi;x]\frac{\delta_L\Delta^m[\chi;y]}{\delta\chi^n(x)}
\right]=0,
\tag{17.2.7}\label{eq:17.2.7}
\end{equation}
respectively. These relations may be made more perspicuous by introducing the quantities
\begin{equation}
\Gamma_N^{(\epsilon)}[\chi]\equiv S_R[\chi]+\epsilon\Gamma_{N,\infty}[\chi,0],
\tag{17.2.8}\label{eq:17.2.8}
\end{equation}
and
\begin{equation}
\Delta_N^{(\epsilon)n}(x)\equiv\Delta^n(x)+\epsilon\mathcal D_N^n(x),
\tag{17.2.9}\label{eq:17.2.9}
\end{equation}
with $\epsilon$ infinitesimal. Then \eqref{eq:17.2.6} (together with the BRST invariance of $S_R$) just says that $\Gamma_N^{(\epsilon)}[\chi]$ is invariant under the transformation
\begin{equation}
\chi^n(x)\longrightarrow\chi^n(x)+\theta\Delta_N^{(\epsilon)n}(x),
\tag{17.2.10}\label{eq:17.2.10}
\end{equation}
while Eq.~\eqref{eq:17.2.7} (together with the nilpotence of the original BRST transformation) tells us that this transformation is \emph{nilpotent}.

We must now consider what form this nilpotent transformation may take. As already mentioned, $\Gamma_{N,\infty}$ consists only of terms of dimensionality four or less, so $\mathcal D_N^n$ and hence $\Delta_N^{(\epsilon)n}(x)$ have at most the dimensionality of the original BRST transformation function $\Delta^n(x)$. Also, $\mathcal D_N^n$ and hence also $\Delta_N^{(\epsilon)n}(x)$ must have the same Lorentz transformation properties and ghost quantum numbers as $\Delta^n(x)$. Hence the most general form of the transformation \eqref{eq:17.2.10} is
\begin{align*}
\psi&\longrightarrow\psi+i\theta\omega^\alpha T_\alpha\psi,\qquad
A_{\alpha\mu}\longrightarrow A_{\alpha\mu}
+\theta\bigl[B_{\alpha\beta}\partial_\mu\omega_\beta
+D_{\alpha\beta\gamma}A_{\beta\mu}\omega_\gamma\bigr],\notag\\
\omega_\alpha&\longrightarrow\omega_\alpha
-\frac12\theta E_{\alpha\beta\gamma}\omega_\beta\omega_\gamma,
\end{align*}
where $T_\alpha$ is some matrix acting on the spinor fields, and $B_{\alpha\beta}$, $D_{\alpha\beta\gamma}$, and $E_{\alpha\beta\gamma}$ are constants, with $E_{\alpha\beta\gamma}$ antisymmetric in $\beta$ and $\gamma$. Also, the transformations of $\omega_\alpha^*$ and $h_\alpha$ are linear, and are therefore unchanged:
\[
\omega_\alpha^*\longrightarrow\omega_\alpha^*-\theta h_\alpha,
\qquad h_\alpha\longrightarrow h_\alpha.
\]

Next we impose the condition of nilpotence. The most important requirement is that $E_{\alpha\beta\gamma}\omega_\beta\omega_\gamma$ should be invariant. This yields the requirement that $E_{\alpha\beta\gamma}E_{\beta\delta\epsilon}\omega_\delta\omega_\epsilon\omega_\gamma$ should vanish, so that the part of $E_{\alpha\beta\gamma}E_{\beta\delta\epsilon}$ that is totally antisymmetric in $\delta$, $\epsilon$, $\gamma$ vanishes:
\[
E_{\alpha\beta\gamma}E_{\beta\delta\epsilon}
+E_{\alpha\beta\epsilon}E_{\beta\gamma\delta}
+E_{\alpha\beta\delta}E_{\beta\epsilon\gamma}=0.
\]
But this just tells us that $E_{\alpha\beta\gamma}$ is the structure constant of some Lie algebra $\mathcal E$. Because $E_{\alpha\beta\gamma}$ goes to the structure constant $C_{\alpha\beta\gamma}$ of the original gauge Lie algebra $\mathcal A$ for $\epsilon\to0$, $\mathcal E$ must be the same as $\mathcal A$, and the structure constants $E_{\alpha\beta\gamma}$ can differ from the original $C_{\alpha\beta\gamma}$ only by a multiplicative factor:\footnote{The requirement of global gauge invariance rules out any non-trivial similarity transformation in the relation between $E_{\alpha\beta\gamma}$ and $C_{\alpha\beta\gamma}$.}
\[
E_{\alpha\beta\gamma}=\mathcal Z C_{\alpha\beta\gamma}.
\]
(This is for simple gauge groups; in the general case we would have a separate factor $\mathcal Z$ for each simple subgroup.)

Next we turn to the condition that the transformation \eqref{eq:17.2.10} be nilpotent when acting on the gauge fields. The requirement that $B_{\alpha\beta}\partial_\mu\omega_\beta+D_{\alpha\beta\gamma}A_{\beta\mu}\omega_\gamma$ be invariant tells us that
\[
D_{\alpha\beta\gamma}D_{\beta\delta\epsilon}
-D_{\alpha\beta\epsilon}D_{\beta\delta\gamma}
=E_{\beta\epsilon\gamma}D_{\alpha\delta\beta}
=\mathcal Z C_{\beta\epsilon\gamma}D_{\alpha\delta\beta}
\]
and
\[
B_{\alpha\beta}E_{\beta\gamma\delta}=D_{\alpha\beta\delta}B_{\beta\gamma}.
\]
The first condition has the unique solution\footnote{The matrix $(D_\gamma)_{\alpha\beta}=D_{\gamma\alpha\beta}/\mathcal Z$ satisfies the commutation relations of the gauge Lie algebra, $[D_\gamma,D_\epsilon]=C_{\beta\epsilon\gamma}D_\beta$. But the only representation of a simple Lie algebra with the same dimensionality and $\mathcal A$ transformation properties as the adjoint representation of $\mathcal A$ is the adjoint representation itself.}
\[
D_{\alpha\beta\gamma}=\mathcal Z C_{\alpha\beta\gamma}.
\]
The second condition tells us that the matrix $B_{\alpha\beta}$ commutes with the adjoint representation of the gauge group, and hence (since we have chosen the structure constants totally antisymmetric) must be proportional to a Kronecker delta, with a coefficient we shall call $\mathcal Z\mathcal N$:
\[
B_{\alpha\beta}=\mathcal Z\mathcal N\delta_{\alpha\beta}.
\]

Finally, the condition that the transformation \eqref{eq:17.2.10} be nilpotent when acting on the fermion fields (if any) requires that $\omega^\alpha T_\alpha\psi$ is invariant. This tells us that
\[
[T_\beta,T_\gamma]=iE_{\alpha\beta\gamma}T_\alpha,
\]
so $T_\alpha$ differs from the generator $t_\alpha$ in the original Lagrangian only by a factor $\mathcal Z$:
\[
T_\alpha=\mathcal Zt_\alpha.
\]
We have thus seen that, apart from the appearance of the new constants $\mathcal Z$ and $\mathcal N$, the transformation \eqref{eq:17.2.10} is just the BRST transformation with which we started:
\begin{align}
\psi&\longrightarrow\psi+i\mathcal Z\theta\omega^\alpha t_\alpha\psi,
&&\tag{17.2.11}\label{eq:17.2.11}\\
A_{\alpha\mu}&\longrightarrow A_{\alpha\mu}
+\mathcal Z\theta\bigl[\mathcal N\partial_\mu\omega_\alpha
+C_{\alpha\beta\gamma}A_{\beta\mu}\omega_\gamma\bigr],
&&\tag{17.2.12}\label{eq:17.2.12}\\
\omega_\alpha&\longrightarrow\omega_\alpha
-\frac12\mathcal Z\theta C_{\alpha\beta\gamma}\omega_\beta\omega_\gamma,
&&\tag{17.2.13}\label{eq:17.2.13}\\
\omega_\alpha^*&\longrightarrow\omega_\alpha^*-\theta h_\alpha,
&&\tag{17.2.14}\label{eq:17.2.14}\\
h_\alpha&\longrightarrow h_\alpha.
&&\tag{17.2.15}\label{eq:17.2.15}
\end{align}
Now we must use this symmetry to constrain the structure of the corrected action \eqref{eq:17.2.8}. Since this contains only the original renormalized action plus the infinite part of the $N$-loop contribution, it must be the integral of a Lagrangian density
\begin{equation}
\Gamma_N^{(\epsilon)}=\int d^4x\,\mathcal L_N^{(\epsilon)}
\tag{17.2.16}\label{eq:17.2.16}
\end{equation}
with $\mathcal L_N^{(\epsilon)}$ a local function of fields and field derivatives of dimensionality (in powers of mass) no greater than 4. Furthermore, as we found in Section 16.4, $\mathcal L_N^{(\epsilon)}$ must be invariant under all the symmetries of the original Lagrangian that act \emph{linearly} on the fields. To identify these symmetries, recall that in generalized $\xi$-gauge, the `new' Lagrangian density in Eq.~\eqref{eq:15.7.6} takes a form given by replacing the term $-(\partial_\mu A_\alpha^\mu)(\partial_\nu A_\alpha^\nu)/2\xi$ in Eq.~\eqref{eq:15.6.16} with the terms $h_\alpha f_\alpha+\tfrac12\xi h_\alpha h_\alpha$ in Eq.~\eqref{eq:15.7.6}:
\begin{align}
\mathcal L_{\mathrm{NEW}}={}&\mathcal L_M
-\frac14F_\alpha{}^{\mu\nu}F_{\alpha\mu\nu}
-\partial_\mu\omega_\alpha^*\partial^\mu\omega_\alpha
\notag\\
&+C_{\alpha\beta\gamma}(\partial_\mu\omega_\alpha^*)A_\gamma{}^\mu\omega_\beta
+h_\alpha\partial_\mu A_\alpha{}^\mu
+\frac12\xi h_\alpha h_\alpha.
\tag{17.2.17}\label{eq:17.2.17}
\end{align}
Inspection of this formula reveals the following linear symmetries:

\noindent(1) \textbf{Lorentz invariance.}

\noindent(2) \textbf{Global gauge invariance}---that is, invariance under the transformations
\begin{align}
\delta\psi_\ell(x)&=i\epsilon^\alpha(t_\alpha)_\ell{}^m\psi_m(x),
&&\tag{17.2.18}\label{eq:17.2.18}\\
\delta A^\beta{}_\mu(x)&=C_{\beta\gamma\alpha}\epsilon^\alpha A^\gamma{}_\mu(x),
&&\tag{17.2.19}\label{eq:17.2.19}\\
\delta\omega_\beta(x)&=C_{\beta\gamma\alpha}\epsilon^\alpha\omega_\gamma(x),
&&\tag{17.2.20}\label{eq:17.2.20}\\
\delta\omega_\beta^*(x)&=C_{\beta\gamma\alpha}\epsilon^\alpha\omega_\gamma^*(x),
&&\tag{17.2.21}\label{eq:17.2.21}\\
\delta h_\beta(x)&=C_{\beta\gamma\alpha}\epsilon^\alpha h_\gamma(x),
&&\tag{17.2.22}\label{eq:17.2.22}
\end{align}
with constant parameters $\epsilon^\alpha$.

\noindent(3) \textbf{Antighost translation invariance}---that is, invariance under the transformation
\begin{equation}
\omega_\alpha^*(x)\longrightarrow\omega_\alpha^*(x)+c_\alpha,
\tag{17.2.23}\label{eq:17.2.23}
\end{equation}
with arbitrary constant parameters $c_\alpha$.

\noindent(4) \textbf{Ghost number conservation}---that is, the conservation of a ghost number equal to $+1$ for $\omega_\alpha$, $-1$ for $\omega_\alpha^*$, and 0 for all other fields.

We shall now proceed to work out the structure of the most general Lagrangian density that is renormalizable, in the sense that it consists only of terms of dimensionality $+4$ or less, that has these linearly acting symmetries, and that is invariant under the modified BRST transformations \eqref{eq:17.2.11}--\eqref{eq:17.2.15}.

From Eq.~\eqref{eq:17.2.17} we may conclude that the fields $A_\alpha^\mu$, $\omega_\alpha$, $\omega_\alpha^*$, and $h_\alpha$ have the dimensionalities $+1$, $+1$, $+1$, and $+2$, respectively (in powers of mass). Note also that ghost number conservation requires that $\omega$ and $\omega^*$ come in pairs, while antighost translation invariance dictates that $\omega^*$ always appears as a derivative. Each pair of $\omega$ and $\partial_\mu\omega^*$ fields adds $+3$ to the dimensionality, so renormalizability rules out any term with more than one such pair. With one such pair we can have at most one more derivative or one additional gauge field, and Lorentz invariance dictates that we must have one or the other. The only renormalizable allowed interactions involving ghost fields are then linear combinations of terms of the form $\partial_\mu\omega_\alpha^*\partial^\mu\omega_\beta$ or $\partial_\mu\omega_\alpha^*A_\gamma^\mu\omega_\beta$.

Next, let us consider the terms that involve the field $h_\alpha$ and possibly other fields but not $\omega$ or $\omega^*$. This field has dimensionality $+2$, so renormalizability and Lorentz invariance allows this field to appear only\footnote{In a theory with scalar fields, we could also have renormalizable terms with $h_\alpha$ multiplied with one or two scalar fields. Such terms cause no trouble, but for brevity they will not be considered here.} multiplied with another $h_\beta$ or $\partial_\mu A_\beta^\mu$ or $A_\beta^\mu A_{\gamma\mu}$.

Finally, the Lagrangian will contain renormalizable terms involving only the matter and gauge fields. We will call the sum of these terms $\mathcal L_{\psi A}$. Putting this all together and using global gauge invariance, the most general renormalizable interaction allowed by the assumed symmetries (aside from BRST invariance) takes the form:
\begin{align}
\mathcal L_N^{(\epsilon)}={}&\mathcal L_{\psi A}
+\frac12\xi' h_\alpha h_\alpha
+c h_\alpha\partial_\mu A_\alpha^\mu
-e_{\alpha\beta\gamma}h_\alpha A_\beta^\mu A_{\gamma\mu}
\notag\\
&-Z_\omega(\partial_\mu\omega_\alpha^*)(\partial^\mu\omega_\alpha)
-d_{\alpha\beta\gamma}(\partial_\mu\omega_\alpha^*)\omega_\beta A_\gamma^\mu,
\tag{17.2.24}\label{eq:17.2.24}
\end{align}
where $\xi'$, $Z_\omega$, $c$, $d_{\alpha\beta\gamma}$, and $e_{\alpha\beta\gamma}$ are unknown constants, with no constraints except obvious symmetry properties such as global gauge invariance and $e_{\alpha\beta\gamma}=e_{\alpha\gamma\beta}$. (As mentioned earlier, we are assuming for simplicity that the gauge group is simple, but the extension to a direct sum of simple and $U(1)$ gauge groups would be trivial; for instance, instead of one term proportional to $h_\alpha h_\alpha$ we would have a sum of such terms, one for each simple subgroup of the gauge group.)

Now we impose BRST invariance. The cancellation of terms in $\delta\mathcal L_N^{(\epsilon)}$ proportional to $\theta\partial_\mu h_\alpha\partial^\mu\omega_\alpha$ tells us that
\begin{equation}
c=\frac{Z_\omega}{\mathcal Z\mathcal N}.
\tag{17.2.25}\label{eq:17.2.25}
\end{equation}
The cancellation of terms in $\delta\mathcal L_N^{(\epsilon)}$ proportional to $\theta\partial_\mu h_\alpha\omega_\beta A_\gamma^\mu$ (or $\theta\partial_\mu\omega_\alpha^*\omega_\beta\partial^\mu\omega_\gamma$) requires that
\begin{equation}
d_{\alpha\beta\gamma}=-(Z_\omega/\mathcal N)C_{\alpha\beta\gamma}.
\tag{17.2.26}\label{eq:17.2.26}
\end{equation}
The terms in $\delta\mathcal L_N^{(\epsilon)}$ proportional to $\theta\partial_\mu\omega_\alpha^*\omega_\beta\omega_\gamma A_\delta^\mu$ automatically then cancel by virtue of the Jacobi identity for the structure constants. The cancellation of terms in $\delta\mathcal L_N^{(\epsilon)}$ proportional to $\theta h_\alpha\partial_\mu\omega_\beta A_\gamma^\mu$ (or $\theta h_\alpha A_\beta^\mu A_{\gamma\mu}\omega_\delta$) yields
\begin{equation}
e_{\alpha\beta\gamma}=0.
\tag{17.2.27}\label{eq:17.2.27}
\end{equation}
Finally, the effect of the infinitesimal transformation \eqref{eq:17.2.10} on matter and gauge fields is the same as a local gauge transformation with gauge parameters $\epsilon_\alpha=\mathcal Z\mathcal N\theta\omega_\alpha$ and gauge couplings renormalized by a factor $1/\mathcal N$ (that is, with $t_\alpha$ and $C_{\alpha\beta\gamma}$ replaced with $\widetilde t_\alpha\equiv t_\alpha/\mathcal N$ and $\widetilde C_{\alpha\beta\gamma}\equiv C_{\alpha\beta\gamma}/\mathcal N$), so the cancellation of terms in $\delta\mathcal L_N^{(\epsilon)}$ with only one factor of $\omega_\alpha$ and no factors $h_\alpha$ or $\omega_\alpha^*$ simply tells us that the Lagrangian $\mathcal L_{\psi A}$ for these fields is \emph{gauge-invariant}, with a renormalized gauge coupling. We conclude then that the most general renormalizable Lagrangian density allowed by our assumed symmetry principles is
\begin{align}
\mathcal L_N^{(\epsilon)}={}&-\frac14Z_A\widetilde F_\alpha{}^{\mu\nu}\widetilde F_{\alpha\mu\nu}
-Z_\psi\bar\psi\gamma^\mu[\partial_\mu-i\widetilde t_\alpha A_{\alpha\mu}]\psi
+\frac12\xi'h_\alpha h_\alpha
\notag\\
&+\frac{Z_\omega}{\mathcal N\mathcal Z}h_\alpha\partial_\mu A_\alpha^\mu
-Z_\omega(\partial_\mu\omega_\alpha^*)(\partial^\mu\omega_\alpha)
+Z_\omega\widetilde C_{\alpha\beta\gamma}(\partial_\mu\omega_\alpha^*)\omega_\beta A_\gamma^\mu,
\tag{17.2.28}\label{eq:17.2.28}
\end{align}
where the tilde on $\widetilde F_\alpha{}^{\mu\nu}$ indicates that the field strength is to be calculated using the renormalized structure constant $\widetilde C_{\alpha\beta\gamma}\equiv C_{\alpha\beta\gamma}/\mathcal N$. But apart from the appearance of a number of new constant coefficients, this is the same Lagrangian with which we started. The new constants in this Lagrangian (including the gauge coupling constant) may be freely shifted by adjusting the $N$th-order terms in the corresponding constants in the original unrenormalized Lagrangian. In particular, we can adjust these terms to make $\Gamma_N^{(\epsilon)}=S_R$, in which case $\Gamma_{N,\infty}=0$, completing the proof.

\begin{center}
$*\ *\ *$
\end{center}

In the above proof we made important use of the accidental invariance of the gauge-fixed Lagrangian \eqref{eq:17.2.17} under the antighost translation transformation \eqref{eq:17.2.23}. This symmetry would not be present for gauge-fixing functionals other than $f_\alpha=\partial_\mu A_\alpha^\mu$, which are not spacetime derivatives. One frequently cited example which preserves Lorentz invariance and global gauge invariance is $f_\alpha=\partial_\mu A_\alpha^\mu+a_{\alpha\beta\gamma}A_\beta^\mu A_{\gamma\mu}$, where $a_{\alpha\beta\gamma}$ is a constant matrix, symmetric in $\beta$ and $\gamma$, which transforms as a tensor under global gauge transformations. (Such constant tensors exist for all $SU(N)$ groups with $N\geq3$.) Another more important example is the background gauge-fixing functional to be introduced in Section 17.4.

The absence of antighost translation invariance does not affect our argument that $\mathcal L_N^{(\epsilon)}$ is a Lorentz- and global gauge-invariant local function of fields and field derivatives of dimensionality no greater than 4, that is invariant under the renormalized BRST transformation \eqref{eq:17.2.11}--\eqref{eq:17.2.15}. But without antighost translation invariance there are new terms in $\mathcal L_N^{(\epsilon)}$ that satisfy these conditions. Since the transformation \eqref{eq:17.2.11}--\eqref{eq:17.2.15} is nilpotent, we can construct such terms as $s'F$, where the transformation \eqref{eq:17.2.11}--\eqref{eq:17.2.15} is written as $\chi^n\to\chi^n+\theta s'\chi^n$, and $F$ is an arbitrary Lorentz- and global gauge-invariant function of ghost number $-1$. One such term is
\[
a_{\alpha\beta\gamma}s'(\omega_\alpha^*A_{\beta\mu}A_\gamma^\mu)
=-a_{\alpha\beta\gamma}\left[
h_\alpha A_{\beta\mu}A_\gamma^\mu
+2\mathcal Z\omega_\alpha^*(\mathcal N\partial_\mu\omega_\beta
+C_{\beta\delta\epsilon}A_{\delta\mu}\omega_\epsilon)A_\gamma^\mu
\right].
\]
This causes no trouble: it is just a renormalized version of the usual ghost and gauge-fixing terms arising from terms $a_{\alpha\beta\gamma}A_\beta^\mu A_{\gamma\mu}$ in the gauge-fixing functional $f_\alpha$. But there is also another possible term of the form
\[
b_{\alpha\beta\gamma}s'(\omega_\alpha^*\omega_\beta^*\omega_\gamma)
=-b_{\alpha\beta\gamma}\left[
2h_\alpha\omega_\beta^*\omega_\gamma
+\frac12\mathcal Z C_{\gamma\delta\epsilon}
\omega_\alpha^*\omega_\beta^*\omega_\delta\omega_\epsilon
\right],
\]
where $b_{\alpha\beta\gamma}$ is a constant, antisymmetric in $\alpha$ and $\beta$, which transforms as a tensor under global gauge transformations. (Such tensors exist for any Lie group; for instance, we could take $b_{\alpha\beta\gamma}$ to be proportional to $C_{\alpha\beta\gamma}$.) But the Faddeev--Popov--De Witt method cannot yield four-ghost interactions in the Lagrangian, so there are no counterterms available to absorb ultraviolet divergences in this term. This is not just a technical obstacle to proving renormalizability; for gauge-fixing functionals like $f_\alpha=\partial_\mu A_\alpha^\mu+a_{\alpha\beta\gamma}A_\beta^\mu A_{\gamma\mu}$, one-loop graphs actually do yield divergences in four-ghost amplitudes that cannot be cancelled by counterterms in the Faddeev--Popov--De Witt Lagrangian.

Aside from avoiding gauge-fixing functionals other than $f_\alpha=\partial_\mu A_\alpha^\mu$, the only solution to this problem seems to be the one mentioned in Section 15.7. We must give up the Faddeev--Popov--De Witt approach, and instead take the action from the beginning as the most general renormalizable function of the gauge, matter, ghost, and auxiliary fields that is invariant under BRST and the other symmetries of the theory. According to the arguments of Section 15.8, the action can be written in the form $I_0+s\Psi$ with $I_0$ ghost-free, and the $S$-matrix is independent of $\Psi$, so we can justify this procedure by quantizing the gauge theory in axial gauge, where ghosts decouple, and then taking $\Psi$ to be anything we like. In particular, we can include $s(\omega^*\omega^*\omega)$ terms in the action that can serve as counterterms to the divergences in four-ghost vertices.
